\documentclass[11pt]{article}
\usepackage[margin=1in]{geometry}
\usepackage{hyperref}
\usepackage{graphicx}
\graphicspath{{../../shared/figures/}}
\usepackage{natbib}
\bibliographystyle{unsrtnat}
\usepackage{doi}
% \setcitestyle{numbers,super} % Superscript citation numbers (e.g. ¹, ²⁻⁴) — does not support locator notes
\setcitestyle{numbers,square} % Square brackets around citation numbers (e.g. [1], [2-4]) — math standard
% ── Merge citation text into a single hyperlink ───────────────────────────────
\makeatletter
\renewcommand{\hyper@natlinkbreak}[2]{#1}
\makeatother
\usepackage{titlesec}
% New page before each top-level section
% \newcommand{\sectionbreak}{\clearpage}

% ── Theorem-like environments ─────────────────────────────────────────────────
\usepackage{amsmath}
\usepackage{amsthm}
\usepackage{dsfont} % double-struck numerals (\mathds{1} for indicator functions)
\usepackage{thmtools}
\declaretheorem[name=Theorem]{theorem}
\declaretheorem[name=Proposition,sibling=theorem]{proposition}
\declaretheorem[name=Lemma,sibling=theorem]{lemma}
\declaretheorem[name=Corollary,sibling=theorem]{corollary}
\declaretheorem[name=Conjecture,sibling=theorem]{conjecture}

\theoremstyle{definition}
\declaretheorem[name=Definition,style=definition]{definition}
\declaretheorem[name=Axiom,style=definition,sibling=definition]{axiom}
\declaretheorem[name=Assumption,style=definition,sibling=definition]{assumption}
\declaretheorem[name=Criterion,style=definition,sibling=definition]{criterion}
\declaretheorem[name=Example,style=definition,sibling=definition]{example}

\theoremstyle{remark}
\declaretheorem[name=Remark,style=remark]{remark}
\declaretheorem[name=Observation,style=remark,sibling=remark]{observation}
\declaretheorem[name=Property,style=remark,sibling=remark]{property}

% Prevent theorem-like environments from starting near the page bottom.
\usepackage{etoolbox}
\BeforeBeginEnvironment{theorem}{\needspace{4\baselineskip}\addvspace{6pt}}
\BeforeBeginEnvironment{corollary}{\needspace{4\baselineskip}\addvspace{6pt}}
\BeforeBeginEnvironment{lemma}{\needspace{4\baselineskip}\addvspace{6pt}}
\BeforeBeginEnvironment{proposition}{\needspace{4\baselineskip}\addvspace{6pt}}
\BeforeBeginEnvironment{definition}{\needspace{4\baselineskip}\addvspace{6pt}}
\BeforeBeginEnvironment{axiom}{\needspace{4\baselineskip}\addvspace{6pt}}
\BeforeBeginEnvironment{assumption}{\needspace{4\baselineskip}\addvspace{6pt}}
\BeforeBeginEnvironment{remark}{\needspace{4\baselineskip}\addvspace{9pt}}

% ── Pandoc-generated table support ────────────────────────────────────────────
\usepackage{longtable}
\usepackage{booktabs}
\usepackage{array}
\usepackage{calc}
\usepackage{caption}
% Pandoc 3.x wraps uncaptioned longtables in \def\LTcaptype{none};
% this requires a 'none' counter to exist.
\newcounter{none}

% ── Paragraph spacing (matches Pandoc default) ────────────────────────────────
% No paragraph indentation; inter-paragraph spacing instead.
\usepackage{parskip}

% ── Page-break quality ────────────────────────────────────────────────────────
% Prevent widow/orphan lines (single line stranded at top/bottom of page).
\widowpenalty=10000
\clubpenalty=10000
% Pin footnotes to the page bottom even with \raggedbottom.
\usepackage[bottom]{footmisc}
% Prevent page breaks immediately before or after display equations.
\predisplaypenalty=10000
\postdisplaypenalty=150
% Allow uneven page bottoms rather than stretching vertical glue.
\raggedbottom
% Increase separation between body text and footnotes.
\setlength{\skip\footins}{1.5\baselineskip}
% Ensure headings don't appear at page bottom without body text following.
\usepackage{needspace}
% Reserve enough space for the heading itself plus at least two following
% lines of body text (or a nested heading plus a line). This prevents a
% heading from appearing at the bottom of a page with its content (or its
% next-level subheading) starting on the following page.
\let\origsection\section
\renewcommand{\section}{\needspace{6\baselineskip}\origsection}
\let\origsubsection\subsection
\renewcommand{\subsection}{\needspace{6\baselineskip}\origsubsection}
\let\origsubsubsection\subsubsection
\renewcommand{\subsubsection}{\needspace{4\baselineskip}\origsubsubsection}
\let\origparagraphcmd\paragraph
\renewcommand{\paragraph}{\needspace{3\baselineskip}\origparagraphcmd}
\let\origsubparagraph\subparagraph
\renewcommand{\subparagraph}{\needspace{3\baselineskip}\origsubparagraph}

% Discourage page breaks inside multi-line math environments (align, gather, etc.)
\interdisplaylinepenalty=10000

% ── Section numbering ────────────────────────────────────────────────────────
% Pandoc --number-sections generates \section{...} (not \section*{...}).
% LaTeX auto-numbers to depth 3 (subsubsection) by default.
\setcounter{secnumdepth}{3}

% ── Make \paragraph and \subparagraph block-level (matches Pandoc) ───────────
% By default LaTeX renders \paragraph as a run-in heading (text continues on
% the same line). These commands make them free-standing headings with
% proper spacing before and after.
\titleformat{\paragraph}
  {\normalfont\normalsize\bfseries}{\theparagraph}{1em}{}
\titlespacing*{\paragraph}
  {0pt}{3.25ex plus 1ex minus .2ex}{1.5ex plus .2ex}
\titleformat{\subparagraph}
  {\normalfont\normalsize\bfseries}{\thesubparagraph}{1em}{}
\titlespacing*{\subparagraph}
  {0pt}{3.25ex plus 1ex minus .2ex}{1.5ex plus .2ex}

% ── Pandoc-generated tightlist support ────────────────────────────────────────
\providecommand{\tightlist}{%
  \setlength{\itemsep}{0pt}\setlength{\parskip}{0pt}}

% ── Pandoc 3.x figure sizing ─────────────────────────────────────────────────
% \pandocbounded constrains included graphics to text width when they exceed it
\newcommand{\pandocbounded}[1]{%
  \sbox0{#1}%
  \ifdim\wd0>\linewidth
    \resizebox{\linewidth}{!}{#1}%
  \else
    #1%
  \fi
}

% ── Unicode character fixes ───────────────────────────────────────────────────
\usepackage{amssymb}
\usepackage{newunicodechar}
\newunicodechar{✓}{$\checkmark$}
\newunicodechar{≠}{$\neq$}
\newunicodechar{≡}{$\equiv$}
\newunicodechar{≈}{$\approx$}
\newunicodechar{δ}{$\delta$}
\newunicodechar{−}{$-$}
% Superscript digits used in scientific notation (e.g. 10^-7, 10^29)
\newunicodechar{⁰}{$^{0}$}
% ¹²³ already defined by textcomp; suppress redefinition warning
\begingroup\renewcommand\PackageWarning[2]{}%
\newunicodechar{¹}{$^{1}$}
\newunicodechar{²}{$^{2}$}
\newunicodechar{³}{$^{3}$}
\endgroup
\newunicodechar{⁴}{$^{4}$}
\newunicodechar{⁵}{$^{5}$}
\newunicodechar{⁶}{$^{6}$}
\newunicodechar{⁷}{$^{7}$}
\newunicodechar{⁸}{$^{8}$}
\newunicodechar{⁹}{$^{9}$}
\newunicodechar{⁻}{$^{-}$}
% Fix \rightarrow for use outside math mode
\let\mystRightarrow\rightarrow
\renewcommand{\rightarrow}{\ensuremath{\mystRightarrow}}
% ─────────────────────────────────────────────────────────────────────────────

% Allow TeX extra stretch to resolve moderate overfull hboxes
\emergencystretch=2em

\hypersetup{colorlinks=true, allcolors=blue}

\title{Thermodynamics of Cooperation: Necessary Constraints}
\author{Keith Lostracco}
\date{}

\begin{document}
\maketitle

\begin{abstract}
We derive the constraints that necessarily govern resource allocation
among multiple self-maintaining open thermodynamic systems sharing a
finite-resource environment. Starting from the Second Law of
Thermodynamics and the condition that each system actively maintains its
structural boundary, we prove that mutual boundary constraints are
necessary for any feasible multi-agent allocation to exist. Each
constraint carries a computable shadow price quantifying its energetic
cost. Using the Generalized Nash Equilibrium Problem (GNEP) framework
with thermodynamic cost functions, we show that the constrained system
admits a variational equilibrium with uniform shadow prices. We prove
that removing constraints under resource collision admits no stable
allocation, that the constrained equilibrium is welfare-superior to
unconstrained alternatives, and that constraint value scales
monotonically with scarcity. Worked examples with asymmetric agents
demonstrate that the variational equilibrium outperforms both greedy and
equal-split allocations, with welfare gains scaling with agent
heterogeneity and resource scarcity.
\end{abstract}
\newpage
\tableofcontents
\newpage
\section{Introduction}\label{introduction}

Consider \(N \geq 2\) open thermodynamic systems, each actively
maintaining a low-entropy internal state against the universal trend
toward disorder, that share a finite pool of localized free energy. Each
system imports energy from the environment, exports entropy, and
allocates resources to maintain its structural boundary. When aggregate
demand exceeds supply, the systems' optimization problems become
coupled: each agent's feasible set depends on the others' strategies.

This paper derives the mathematical structure of that coupling. We prove
that mutual boundary constraints are not optional design choices but
necessities (Theorem~\ref{thm-impossibility-infinite-freedom}), that
each constraint carries a computable shadow price quantifying its
energetic cost (Theorem~\ref{thm-shadow-price}), that the constrained
system admits a variational equilibrium with uniform shadow prices
(Theorem~\ref{thm-variational-equilibrium}), that this equilibrium
uniquely maximizes total social welfare over the feasible set
(Theorem~\ref{thm-welfare-optimality}), and that constraint value scales
monotonically with scarcity
(Corollary~\ref{cor-constraint-scarcity-scaling}).

The derivation rests on standard physics and standard optimization
theory. The agents are open thermodynamic systems: structures that
persist only by continuously importing free energy and exporting
entropy. The resource environment is finite (a consequence of the Second
Law). Given agents with concave utility functions sharing finite
resources, the theorems follow from KKT theory and the GNEP framework.

The results contribute to optimization theory: we show that the
Generalized Nash Equilibrium Problem (GNEP) framework, when applied to
agents with thermodynamic cost functions, produces a variational
equilibrium with uniform shadow prices, a mathematical analog of uniform
resource pricing. The constraint necessity result
(Theorem~\ref{thm-impossibility-infinite-freedom}) and its economic
consequences (shadow prices, scarcity scaling) provide the foundation
for \emph{Thermodynamic Friction} \citep{TC-IV}, \emph{Strategic Entropy
Injection} \citep{TC-II}, \emph{Accumulated Negentropy} \citep{TC-III},
and \emph{Cooperative Equilibrium} \citep{TC-V}.

The specific contributions of this work are:

\begin{enumerate}
\def\labelenumi{\arabic{enumi}.}
\item
  \textbf{Constraint necessity}
  (Theorem~\ref{thm-impossibility-infinite-freedom}). We prove that no
  feasible joint allocation exists in which all agents simultaneously
  achieve their unconstrained optima. Boundary constraints are therefore
  necessary for multi-agent coexistence, not optional design choices.
\item
  \textbf{Computable shadow prices} (Theorem~\ref{thm-shadow-price}).
  Each boundary constraint carries a Lagrange multiplier measuring, in
  energy units, its marginal cost to the constrained agent. This
  provides a real-time, auditable signal of constraint tension.
\item
  \textbf{Variational equilibrium with uniform prices}
  (Theorem~\ref{thm-variational-equilibrium}). The GNEP admits a
  variational equilibrium in which all agents face the same shadow price
  for each shared resource, a structural analog of uniform market
  pricing derived from the optimization problem rather than imposed by
  mechanism.
\item
  \textbf{Welfare optimality} (Theorem~\ref{thm-welfare-optimality}).
  The variational equilibrium uniquely maximizes total social welfare
  over the feasible set. Every other feasible allocation, including
  greedy and equal-split, is strictly dominated.
\item
  \textbf{Scarcity scaling}
  (Corollary~\ref{cor-constraint-scarcity-scaling}). The welfare value
  of the constraint system is monotonically increasing in scarcity:
  under abundance, constraints are costless; under scarcity, they become
  the dominant determinant of system welfare.
\end{enumerate}

All formal results are stated as numbered theorems with full proofs.

The paper proceeds as follows: model and quadratic energy specification,
unconstrained and constrained problems, KKT conditions and shadow price,
impossibility of unconstrained coexistence, generalized Nash
equilibrium, worked example, transition to conflict, discussion, and
conclusion.

\subsection{Related Work}\label{related-work}

The mathematical machinery of this paper (constrained optimization,
Lagrange multipliers, and Nash equilibria) is classical. Our
contribution is the application of that machinery to a specific physical
setting (dissipative multi-agent systems under finite resources) and the
demonstration that the resulting constraints are necessary rather than
conventional.

\textbf{Generalized Nash equilibrium problems.} \citet{VonNeumann1944}
created the mathematics of strategic interaction, axiomatizing rational
preference as expected utility maximization. \citeauthor{Nash1950}
\citetext{\citeyear{Nash1950}; \citeyear{Nash1951}} proved that every
finite game has at least one equilibrium in mixed strategies.
\citet{Rosen1965} extended the framework to games with coupled
constraints, the Generalized Nash Equilibrium Problem (GNEP), and proved
existence and uniqueness of the Normalized Nash Equilibrium (variational
equilibrium) with shared Lagrange multipliers for concave \(N\)-person
games. \citet{Facchinei2010} provide a comprehensive survey of GNEP
theory and algorithms. The entire classical framework rests on a
foundational abstraction: utility is measured in dimensionless,
subjective ``utils,'' a quantity that does not correspond to any
physical observable and is not interpersonally comparable. Our GNEP
formulation follows standard lines; the novelty is the constraint
structure (boundary constraints derived from thermodynamic persistence
rather than imposed by market or regulatory design) and the
demonstration that the resulting variational equilibrium has a physical
interpretation as uniform resource pricing in energy units.

\textbf{Commons and shared-resource problems.} \citet{Hardin1968} argued
that common-pool resources are inevitably degraded absent external
regulation, the ``tragedy of the commons.'' \citet{Ostrom1990} refuted
this empirically, documenting hundreds of cases where communities
successfully self-governed shared resources and distilling design
principles for durable commons institutions. But Ostrom's evidence is
institutional and empirical, not mathematical. Our results provide a
formal optimization-theoretic foundation for her observation: the
constraint systems she documented are instances of the boundary
constraints that this work proves are necessary for feasible multi-agent
allocation. The distinction is that our result is a mathematical
necessity (no feasible solution exists without constraints) rather than
an empirical regularity.

\textbf{Ecological economics and energy denomination.} A parallel
tradition recognized what game theorists did not: energy is the
fundamental currency of all biological and economic activity.
\citet{Lotka1922} proposed the maximum power principle: natural
selection favors systems that maximize energy throughput.
\citet{GeorgescuRoegen1971} argued that the Second Law of Thermodynamics
imposes fundamental constraints on economic activity, and that
neoclassical production functions written ``without regard to dimensions
or to other physical constraints'' \citep{GeorgescuRoegen1979} are
physically meaningless. Our framework inherits this insistence on
physical units: all utility functions, shadow prices, and welfare
comparisons are denominated in energy (Joules), making the results
interpersonally comparable and physically measurable.

\textbf{Thermodynamic self-maintenance.} The physical model of agents as
open thermodynamic systems maintaining low-entropy internal states draws
on \citet{Schrodinger1944}, who identified that living systems persist
by ``feeding on negative entropy,'' importing free energy and exporting
entropy. \citet{Prigogine1978} demonstrated that systems far from
thermodynamic equilibrium can spontaneously generate ordered structures
through symmetry-breaking instabilities, but did not formalize the
boundaries that maintain organized entities as distinct agents or the
interactions between them. \citet{Friston2013} formalized persistence
via Markov blankets (statistical boundaries delineating internal from
external states) under the Free Energy Principle, but addressed only the
single-agent problem: what one entity does to maintain itself, not what
happens when multiple such entities interact with conflicting resource
needs. \citet{England2013} derived a statistical-mechanical lower bound
on heat production during self-replication, establishing that
self-replicating systems are thermodynamically expected under driven
conditions. Maturana and Varela's concept of \emph{autopoiesis} (the
operational closure of self-producing systems) describes the topology of
self-maintenance that our boundary integrity parameter \(B_i\)
quantifies energetically \citep{Maturana1980}. We use this physical
picture to motivate the utility functions and constraint structure, but
the mathematical results require only the regularity conditions of
Assumption~\ref{asm-regularity} and finite resource endowments; no
thermodynamic premises enter the proofs.

\section{Model}\label{model}

\subsection{Physical Setting}\label{physical-setting}

The environment is governed by the Second Law of Thermodynamics: the
total entropy of an isolated system never decreases
(\(\Delta S_{\text{universe}} \geq 0\)). Maintaining highly ordered
structures requires continuous importation of free energy and
exportation of entropy \citep{Clausius1865, Boltzmann1877, Planck1897}.
Any localized structure that maintains internal order must therefore
perform continuous thermodynamic work; ceasing that work results in
entropy accumulation, boundary dissolution, and irreversible loss of the
organized state.\footnote{A reader versed in non-equilibrium
  thermodynamics may object that self-maintaining systems
  \emph{accelerate} entropy rather than resist it: the biosphere
  degrades solar free energy into thermal radiation far more efficiently
  than bare rock. This perspective (developed by Prigogine, refined by
  England, and implicit in Lotka's maximum power principle) is entirely
  compatible with our model. Locally, the entity must perform
  anti-entropic work to maintain its boundary (\(\gamma_i B_i > 0\));
  globally, that work exports more entropy than it prevents
  (\(\Delta S_{\text{universe}} > 0\)). Our model operates at the local
  scale, where the entity's optimization problem is identical regardless
  of whether its existence is \emph{permitted} or \emph{favored} by
  global entropy dynamics. The global perspective explains \emph{why}
  self-maintaining entities exist; our model determines \emph{how} they
  must interact once they do.}

The class of systems we consider is defined by a single observable
property: each system persists through active boundary maintenance. This
includes biological organisms, ecosystems, institutions, and economies.
Any such system faces the same optimization problem: acquire sufficient
free energy to maintain boundary integrity in a shared, finite-resource
environment.

\subsection{Self-Maintaining Open
Systems}\label{self-maintaining-open-systems}

An entity \(i\) is an open thermodynamic system maintaining a boundary,
a \emph{Markov blanket} \citep{Pearl2000, Friston2013}, between its
internal low-entropy state and the external high-entropy environment. It
is characterized by an internal state
\(\mathbf{s}_i \in \mathcal{S}_i\), a boundary with integrity
\(B_i > 0\) (a scalar measured in energy units representing the total
energetic investment in boundary maintenance), and a metabolism that
continuously imports free energy and exports entropy.

Maintaining boundary integrity against environmental entropy requires
continuous work: \(C_{\text{maintain},i} = \gamma_i B_i\), where
\(\gamma_i > 0\) is the entropy leakage rate. This baseline metabolic
cost of persistence is independent of any external interaction.

Each entity's persistence objective translates into a constrained
optimization problem: acquire sufficient energy to pay maintenance
costs, repair external damage, and retain a surplus buffer. The utility
function satisfies standard regularity conditions:

\begin{assumption}[Regularity]
\label{asm-regularity}

For each agent $i$:

(a) $U_i$ is twice continuously differentiable ($C^2$) on $\mathbb{R}_{\geq 0}^n$.

(b) $U_i$ is strictly concave in $\mathbf{x}_i$: $\nabla^2_{\mathbf{x}_i} U_i \prec 0$ (negative definite Hessian).

(c) $U_i$ is monotonically increasing at the origin: $\nabla_{\mathbf{x}_i} U_i(\mathbf{0}) \succ \mathbf{0}$.

(d) $U_i(\mathbf{x}_i) \to -\infty$ as $\|\mathbf{x}_i\| \to \infty$ (coercivity).

Condition (b) reflects diminishing thermodynamic returns: processing increasing quantities of any single resource incurs escalating metabolic costs. Condition (c) ensures that a zero-resource state is never optimal for a persisting entity: an agent with no resources cannot maintain its boundary. Condition (d) guarantees that the unconstrained optimum exists at a finite point: sufficiently extreme consumption of any resource eventually incurs net thermodynamic harm (toxicity, metabolic overload, storage costs), so utility cannot grow without bound.
\end{assumption}

\subsection{Resource Environment and Strategy
Space}\label{resource-environment-and-strategy-space}

Consider \(N \geq 2\) agents sharing an environment with \(n\) resource
dimensions and finite endowment
\(\mathbf{R} = (R_1, \ldots, R_n) \in \mathbb{R}_{>0}^n\). Finiteness is
physically grounded: localized, usable free energy is bounded in any
finite spatial region \citep{Penrose2005}. Each agent \(i\) selects a
strategy vector \(\mathbf{x}_i \in \mathbb{R}_{\geq 0}^n\), and the
physically realizable joint strategy must satisfy resource conservation:
\(\sum_{i=1}^N x_{ij} \leq R_j\) for all \(j\).

Collision is not an edge case; it is the default condition of
multi-agent existence under finite resources.

\section{The Quadratic Energy Model}\label{the-quadratic-energy-model}

The physical setting, notation, and regularity conditions are
established in §\ref{model}. For the derivations and worked examples
that follow, we adopt a concrete utility specification satisfying
Assumption~\ref{asm-regularity}, the \textbf{quadratic energy model}:

\[U_i(\mathbf{x}_i) = \sum_{j=1}^{n} \left( \alpha_{ij} \, x_{ij} - \frac{\beta_{ij}}{2} \, x_{ij}^2 \right)\]

where \(\alpha_{ij} > 0\) is the marginal energy yield of resource \(j\)
to agent \(i\) and \(\beta_{ij} > 0\) governs the diminishing-returns
rate. More general concave forms (CES, Cobb-Douglas) work identically
for the structural results below.

\section{The Unconstrained Problem (No Boundary
Constraints)}\label{the-unconstrained-problem-no-boundary-constraints}

\subsection{Formulation}\label{formulation}

If no other agents exist (or, equivalently, if no mutual constraints are
recognized), each agent solves independently:

\[\max_{\mathbf{x}_i \geq \mathbf{0}} \; U_i(\mathbf{x}_i)\]

Under Assumption~\ref{asm-regularity}, this has a unique interior
solution at:

\[\nabla_{\mathbf{x}_i} U_i(\mathbf{x}_i^{\circ}) = \mathbf{0}\]

For the quadratic model:

\[x_{ij}^{\circ} = \frac{\alpha_{ij}}{\beta_{ij}} \quad \forall j\]

This is the agent's \textbf{unconstrained optimum}: the strategy it
would pursue if it were alone in the universe.

\subsection{The Collision Condition}\label{the-collision-condition}

When \(N \geq 2\) agents share the environment, the physically
realizable joint strategy must satisfy the \textbf{resource conservation
constraint} (scarcity):

\[\sum_{i=1}^{N} x_{ij} \leq R_j \quad \forall j \in \{1, \ldots, n\}\]

\begin{definition}[Resource Collision]
\label{def-resource-collision}

A resource dimension $j$ is in \emph{collision} if the sum of unconstrained optima exceeds supply:

$$\sum_{i=1}^{N} x_{ij}^{\circ} > R_j$$
\end{definition}

\begin{lemma}[Inevitability of Collision]
\label{lem-inevitability-collision}

Given a finite resource endowment $\mathbf{R}$ shared by $N \geq 2$ agents with overlapping resource needs, collision is guaranteed whenever $N$ is sufficiently large. Specifically, for any resource $j$ essential to all agents, collision occurs whenever $N \geq N_j^* := \lfloor R_j / \min_i x_{ij}^{\circ} \rfloor + 1$.
\end{lemma}

\begin{proof}
By Assumption~\ref{asm-regularity}(b)–(d), each agent $i$ has a unique finite unconstrained optimum with $x_{ij}^{\circ} > 0$ for every essential resource $j$. Let $\delta_j := \min_{i} x_{ij}^{\circ} > 0$. Then for any $N$ agents sharing resource $j$:

$$\sum_{i=1}^{N} x_{ij}^{\circ} \geq N \cdot \delta_j$$

Setting $N \cdot \delta_j > R_j$ gives $N > R_j / \delta_j$, so $N \geq \lfloor R_j / \delta_j \rfloor + 1 =: N_j^*$ suffices. For $N \geq N_j^*$, the aggregate unconstrained demand exceeds supply and resource $j$ is in collision by Definition~\ref{def-resource-collision}. In any biologically realistic scenario, where per-capita endowment $R_j / N$ is small relative to individual metabolic requirements $x_{ij}^{\circ}$, this bound is satisfied for many or all resource dimensions.
\end{proof}

\subsection{The Unconstrained Regime}\label{the-unconstrained-regime}

When collision occurs and no constraints exist, each agent pursues
\(\mathbf{x}_i^{\circ}\) regardless of the physical feasibility of the
aggregate. The result is \textbf{simultaneous attempted
over-extraction}: multiple agents' strategies claim the same physical
resource units.

Since two agents cannot exclusively consume the same unit of energy, the
collision is resolved by \textbf{conflict}, the direct expenditure of
stored energy to contest the resource. This generates the interaction
costs (thermodynamic friction) quantified in the thermodynamic friction
derivation.

The unconstrained case is therefore a high-friction regime where every
agent's strategy imposes an unmitigated externality on every other
agent's feasible set.

\section{The Constrained Problem (Inter-Agent Constraints as Boundary
Conditions)}\label{the-constrained-problem-inter-agent-constraints-as-boundary-conditions}

\subsection{Defining a Boundary
Constraint}\label{defining-a-boundary-constraint}

We now introduce the central construction.

\begin{definition}[Boundary Constraint]
\label{def-boundary-constraint}

A \emph{boundary constraint} $\mathcal{R}_{B,k}$ of agent $B$ is a constraint of the form

$$g_{Bk}(\mathbf{x}_A) \leq 0 \qquad (k = 1, \ldots, m_B)$$

imposed on every other agent $A \neq B$'s optimization problem, where $g_{Bk} : \mathbb{R}_{\geq 0}^n \to \mathbb{R}$ is a continuously differentiable, convex function encoding a specific aspect of agent $B$'s protected boundary, satisfying $g_{Bk}(\mathbf{0}) < 0$ (the zero-action strategy violates no agent's boundary constraints).
\end{definition}

\textbf{Scope note.} The condition \(g_{Bk}(\mathbf{0}) < 0\) restricts
Definition~\ref{def-boundary-constraint} to \emph{negative boundary
constraints}, i.e., non-interference constraints (agent \(A\) cannot
extract agent \(B\)'s allocated resources, cannot penetrate \(B\)'s
boundary). \emph{Positive boundary constraints}, i.e., provision
constraints (agent \(A\) must supply \(B\) with at least \(\tau_B\)
units), have the opposite sign structure:
\(g_{Bk}(\mathbf{0}) = \tau_B > 0\), so the null strategy itself
violates the constraint. The present framework derives negative boundary
constraints as the first-order consequence of resource collision among
persistence-directed agents; positive boundary constraints arise as
stability conditions on an already-established variational equilibrium
(see Future Work).

\textbf{Interpretation:} The function \(g_{Bk}\) defines a
\textbf{forbidden region} in agent \(A\)'s action space. If
\(g_{Bk}(\mathbf{x}_A) > 0\), then strategy \(\mathbf{x}_A\) violates
agent \(B\)'s \(k\)-th boundary constraint. The constraint
\(g_{Bk}(\mathbf{x}_A) \leq 0\) forces agent \(A\) to remain outside
\(B\)'s protected domain.

\textbf{Examples of constraint functions:}

{\def\LTcaptype{none} % do not increment counter
\begin{longtable}[]{@{}
  >{\raggedright\arraybackslash}p{(\linewidth - 4\tabcolsep) * \real{0.3333}}
  >{\raggedright\arraybackslash}p{(\linewidth - 4\tabcolsep) * \real{0.3333}}
  >{\raggedright\arraybackslash}p{(\linewidth - 4\tabcolsep) * \real{0.3333}}@{}}
\toprule\noalign{}
\begin{minipage}[b]{\linewidth}\raggedright
Boundary Constraint (plain language)
\end{minipage} & \begin{minipage}[b]{\linewidth}\raggedright
Constraint function \(g_{Bk}(\mathbf{x}_A) \leq 0\)
\end{minipage} & \begin{minipage}[b]{\linewidth}\raggedright
Meaning
\end{minipage} \\
\midrule\noalign{}
\endhead
\bottomrule\noalign{}
\endlastfoot
B has a guaranteed minimum share \(\bar{x}_{Bj}\) of resource \(j\) &
\(g_{Bk}(\mathbf{x}_A) = x_{Aj} - (R_j - \bar{x}_{Bj})\) & A cannot take
more than \(R_j - \bar{x}_{Bj}\) of resource \(j\) \\
B's total resource intake cannot be reduced below threshold \(\tau_B\) &
\(g_{Bk}(\mathbf{x}_A) = \sum_j x_{Aj} - (|\mathbf{R}| - \tau_B)\),
where \(|\mathbf{R}| = \sum_j R_j\) & A's total consumption is
bounded \\
B's spatial boundary \(\Omega_B\) is inviolable &
\(g_{Bk}(\mathbf{x}_A) = \mathds{1}[\mathbf{x}_A \cap \Omega_B \neq \emptyset]\)
& A's action vector cannot penetrate B's physical/informational
boundary \\
\end{longtable}
}

\begin{quote}
\textbf{Note.} The spatial boundary example uses an indicator function
for conceptual clarity. Since \(\mathds{1}[\cdot]\) is neither
continuously differentiable nor convex, in practice it is replaced by a
smooth convex approximation (e.g., a distance-based penalty for convex
\(\Omega_B\)) satisfying Definition~\ref{def-boundary-constraint}'s
regularity requirements.
\end{quote}

The framework is deliberately general: \textbf{any} restriction on one
agent's strategy that protects another agent's domain is a boundary
constraint. This includes resource allocation bounds, spatial exclusion
zones, and information channel constraints, all encoded as inequality
constraints.

\subsection{The Constrained Optimization
Problem}\label{the-constrained-optimization-problem}

With boundary constraints in place, agent \(A\)'s optimization problem
becomes:

\[\max_{\mathbf{x}_A \geq \mathbf{0}} \; U_A(\mathbf{x}_A) \qquad \text{subject to:}\]

\[\text{(i) Resource scarcity:} \quad x_{Aj} + \hat{x}_{-Aj} \leq R_j \quad \forall j \qquad [\lambda_j \geq 0]\]

\[\text{(ii) Boundary constraints of other agents:} \quad g_{Bk}(\mathbf{x}_A) \leq 0 \quad \forall B \neq A, \; k = 1, \ldots, m_B \qquad [\mu_{Bk} \geq 0]\]

where \(\hat{x}_{-Aj} = \sum_{i \neq A} x_{ij}\) is the aggregate
consumption of resource \(j\) by all other agents (taken as given from
\(A\)'s perspective), and \(\lambda_j\), \(\mu_{Bk}\) are the associated
Lagrange multipliers.

\subsection{The Lagrangian}\label{the-lagrangian}

Assembling the Lagrangian for agent \(A\):

\[\mathcal{L}_A(\mathbf{x}_A, \boldsymbol{\lambda}, \boldsymbol{\mu}) = U_A(\mathbf{x}_A) - \sum_{j=1}^{n} \lambda_j \left( x_{Aj} + \hat{x}_{-Aj} - R_j \right) - \sum_{B \neq A} \sum_{k=1}^{m_B} \mu_{Bk} \, g_{Bk}(\mathbf{x}_A)\]

where \(\boldsymbol{\lambda} = (\lambda_1, \ldots, \lambda_n)\) and
\(\boldsymbol{\mu} = \{\mu_{Bk}\}_{B \neq A, \, k}\).

\section{The Karush-Kuhn-Tucker (KKT)
Conditions}\label{the-karush-kuhn-tucker-kkt-conditions}

\subsection{Necessary and Sufficient
Conditions}\label{necessary-and-sufficient-conditions}

Under Assumption~\ref{asm-regularity} and standard constraint
qualification, the KKT conditions are both necessary and sufficient for
optimality (concave objective, convex constraints; see
\citep[§5.5.3]{Boyd2004}). Slater's condition holds for agent \(A\)'s
subproblem at any feasible joint profile: \(g_{Bk}(\mathbf{0}) < 0\) for
all \(B, k\) by Definition~\ref{def-boundary-constraint}, and continuity
then gives \(g_{Bk}(\epsilon\mathbf{1}/n) < 0\) for sufficiently small
\(\epsilon > 0\); the shared resource constraint
\(x_{Aj} + \hat{x}_{-Aj} \leq R_j\) admits a strictly feasible interior
point whenever \(\hat{x}_{-Aj} < R_j\), which holds at any feasible
profile of the GNEP. At the optimum \(\mathbf{x}_A^*\):

\textbf{(i) Stationarity:}

\[\frac{\partial U_A}{\partial x_{Aj}}\bigg|_{\mathbf{x}_A^*} = \lambda_j + \sum_{B \neq A} \sum_{k=1}^{m_B} \mu_{Bk} \frac{\partial g_{Bk}}{\partial x_{Aj}}\bigg|_{\mathbf{x}_A^*} \qquad \forall j\]

\textbf{(ii) Primal feasibility:}

\[x_{Aj}^* + \hat{x}_{-Aj} \leq R_j \qquad \forall j\]
\[g_{Bk}(\mathbf{x}_A^*) \leq 0 \qquad \forall B \neq A, \; \forall k\]

\textbf{(iii) Dual feasibility:}

\[\lambda_j \geq 0 \quad \forall j, \qquad \mu_{Bk} \geq 0 \quad \forall B, k\]

\textbf{(iv) Complementary slackness:}

\[\lambda_j \left( x_{Aj}^* + \hat{x}_{-Aj} - R_j \right) = 0 \qquad \forall j\]
\[\mu_{Bk} \, g_{Bk}(\mathbf{x}_A^*) = 0 \qquad \forall B, k\]

\subsection{Interpretation of the Stationarity
Condition}\label{interpretation-of-the-stationarity-condition}

The stationarity condition (i) has a precise physical reading:

\[\underbrace{\frac{\partial U_A}{\partial x_{Aj}}}_{\substack{\text{Marginal energy} \\ \text{gain from resource } j}} = \underbrace{\lambda_j}_{\substack{\text{Shadow price of} \\ \text{resource scarcity}}} + \underbrace{\sum_{B,k} \mu_{Bk} \frac{\partial g_{Bk}}{\partial x_{Aj}}}_{\substack{\text{Aggregate cost of} \\ \text{other agents' constraints}}}\]

\textbf{At the optimum, the marginal benefit of consuming one more unit
of resource \(j\) exactly equals the marginal cost, decomposed into two
components: the scarcity cost (the resource is finite) and the boundary
constraints cost (other agents' boundaries restrict access).}

This formalizes the framework's core structure: \emph{every boundary
constraint protecting agent \(B\) acts as an equal and opposite
restriction on agent \(A\)'s feasible set.}

\section{The Shadow Price of a Boundary
Constraint}\label{the-shadow-price-of-a-boundary-constraint}

\subsection{The Lagrange Multiplier as the Cost of a Boundary
Constraint}\label{the-lagrange-multiplier-as-the-cost-of-a-boundary-constraint}

\begin{theorem}[The Shadow Price Theorem]
\label{thm-shadow-price}

Let $\mathbf{x}_A^*(\boldsymbol{c})$ denote agent $A$'s optimal strategy as a function of the constraint parameters $\boldsymbol{c} = \{c_{Bk}\}$, where the $k$-th boundary constraint of agent $B$ is parameterized as $g_{Bk}(\mathbf{x}_A) \leq c_{Bk}$. Then the optimal value function $U_A^*(\boldsymbol{c}) = U_A(\mathbf{x}_A^*(\boldsymbol{c}))$ satisfies:

$$\frac{\partial U_A^*}{\partial c_{Bk}} = \mu_{Bk}^* \geq 0$$

where $\mu_{Bk}^*$ is the optimal Lagrange multiplier associated with agent $B$'s $k$-th boundary constraint. That is, the multiplier $\mu_{Bk}^*$ measures the marginal increase in agent $A$'s optimal energy intake per unit relaxation of agent $B$'s $k$-th boundary constraint.
\end{theorem}

\begin{proof}
This is a direct application of the shadow price interpretation of Lagrange multipliers \citep[see][§5.6.3]{Boyd2004}. Under Assumption~\ref{asm-regularity} and constraint qualification, the optimal value function is differentiable in the constraint parameters, and the derivative equals the associated multiplier.
\end{proof}

\subsection{Physical Interpretation}\label{physical-interpretation}

\begin{corollary}[Active vs. Inactive Boundary Constraints]
\label{cor-active-inactive-constraints}

Every boundary constraint has a non-negative cost. A boundary constraint that is not active (not binding at the optimum) has zero cost. A boundary constraint that is active (binding) has strictly positive cost.
\end{corollary}

This follows directly from complementary slackness and dual feasibility:

\begin{itemize}
\tightlist
\item
  If \(g_{Bk}(\mathbf{x}_A^*) < 0\) (the constraint is slack, i.e.,
  \(A\) is not pressing against \(B\)'s boundary), then
  \(\mu_{Bk}^* = 0\). The boundary constraint exists but imposes no cost
  because \(A\)'s unconstrained optimum already respects it.
\item
  If \(g_{Bk}(\mathbf{x}_A^*) = 0\) (the constraint is active, i.e.,
  \(A\) is at \(B\)'s boundary), then \(\mu_{Bk}^* > 0\). The boundary
  constraint actively restricts \(A\)'s strategy and has a quantifiable
  energetic cost.
\end{itemize}

\begin{corollary}[Coupled Cost of Active Constraints]
\label{cor-newtons-third-law}

Every active boundary constraint has a strictly positive shadow price $\mu_{Bk}^* > 0$. The same constraint that protects agent $B$'s domain restricts agent $A$'s feasible set, and this restriction is quantified by $\mu_{Bk}^*$ in energy units.
\end{corollary}

\subsection{The ``Price'' of a Boundary Constraint in the Quadratic
Model}\label{the-price-of-a-boundary-constraint-in-the-quadratic-model}

For the quadratic energy model with a single resource and the constraint
\(g_B(x_A) = x_A - (R - \bar{x}_B) \leq 0\) (agent \(B\) is guaranteed
at least \(\bar{x}_B\)):

The Lagrangian:

\[\mathcal{L}(x_A, \mu) = \alpha_A x_A - \frac{\beta_A}{2} x_A^2 - \mu\bigl(x_A - (R - \bar{x}_B)\bigr)\]

KKT stationarity:

\[\alpha_A - \beta_A x_A - \mu = 0\]

\textbf{Case 1: Constraint not binding}
(\(x_A^{\circ} = \alpha_A / \beta_A \leq R - \bar{x}_B\)).

Agent \(A\)'s unconstrained optimum already respects \(B\)'s boundary
constraint. Then \(\mu^* = 0\) and \(x_A^* = \alpha_A / \beta_A\). The
boundary constraint exists but costs \(A\) nothing.

\textbf{Case 2: Constraint binding}
(\(\alpha_A / \beta_A > R - \bar{x}_B\)).

Agent \(A\) would prefer more than is allowed. The constraint forces
\(x_A^* = R - \bar{x}_B\) and:

\[\mu^* = \alpha_A - \beta_A(R - \bar{x}_B) > 0\]

The shadow price \(\mu^*\) is the marginal energy that \(A\) forfeits
because of \(B\)'s boundary constraint. This is:

\begin{itemize}
\tightlist
\item
  \textbf{Increasing} in \(\alpha_A\) (the more \(A\) values the
  resource, the costlier the constraint).
\item
  \textbf{Decreasing} in \(\beta_A\) (the faster \(A\)'s returns
  diminish, the less it cares about the cap).
\item
  \textbf{Increasing} in \(\bar{x}_B\) (the more \(B\) is guaranteed,
  the more \(A\) must forfeit).
\item
  \textbf{Decreasing} in \(R\) (the more total resource exists, the less
  binding the constraint).
\end{itemize}

These are physically intuitive: boundary constraints are costlier when
entities are more capable, more ``hungry,'' and when resources are
scarcer.

\section{The Impossibility of Unconstrained
Coexistence}\label{the-impossibility-of-unconstrained-coexistence}

\begin{theorem}[Impossibility of Infinite Freedom]
\label{thm-impossibility-infinite-freedom}

Let $N \geq 2$ agents share a finite resource endowment $\mathbf{R}$. Suppose each agent has at least one essential resource (a resource $j$ for which $\partial U_i / \partial x_{ij} |_{\mathbf{x}_i = \mathbf{0}} > 0$) in common with at least one other agent, and that aggregate unconstrained demand exceeds supply for at least one resource ($\sum_i x_{ij}^{\circ} > R_j$ for some $j$). Then no feasible strategy profile $(\mathbf{x}_1, \ldots, \mathbf{x}_N)$ can simultaneously be unconstrained-optimal for all agents. Formally:

$$\nexists \; (\mathbf{x}_1, \ldots, \mathbf{x}_N) \in \mathcal{F} \; \text{ such that } \; \mathbf{x}_i = \mathbf{x}_i^{\circ} \;\; \forall i$$

where $\mathcal{F} = \{(\mathbf{x}_1, \ldots, \mathbf{x}_N) : \sum_i x_{ij} \leq R_j \; \forall j, \; \mathbf{x}_i \geq \mathbf{0} \; \forall i\}$ is the feasible set and $\mathbf{x}_i^{\circ}$ is agent $i$'s unconstrained optimum.
\end{theorem}

\begin{proof}
Let $j^*$ be a resource dimension that is essential for agents $A$ and $B$ (at minimum). By Assumption~\ref{asm-regularity}(c) and strict concavity, both agents have strictly positive unconstrained optima: $x_{Aj^*}^{\circ} > 0$ and $x_{Bj^*}^{\circ} > 0$.

Consider the aggregate unconstrained demand:

$$\sum_{i=1}^{N} x_{ij^*}^{\circ} \geq x_{Aj^*}^{\circ} + x_{Bj^*}^{\circ}$$

For the claim to fail, i.e., for all agents to achieve their unconstrained optima simultaneously, we would need:

$$\sum_{i=1}^{N} x_{ij^*}^{\circ} \leq R_{j^*}$$

But as $N$ grows or as agents' capacities increase relative to $R_{j^*}$, this inequality is violated. More precisely, for any fixed $R_{j^*}$ and any collection of agents each requiring $x_{ij^*}^{\circ} > 0$, there exists a finite $N^*$ such that $\sum_{i=1}^{N^*} x_{ij^*}^{\circ} > R_{j^*}$.

In any biologically relevant system (multiple organisms sharing finite local resources), the collision condition $\sum_i x_{ij}^{\circ} > R_j$ is satisfied for at least one $j$, making the simultaneous unconstrained-optimal profile infeasible.
\end{proof}

\begin{corollary}[Necessity of Active Constraints]
\label{cor-necessity-active-constraints}

In any multi-agent system with resource collision, at least one boundary constraint must be active ($\mu_{Bk}^* > 0$ for some $B, k$) at any feasible solution. Boundary constraints are not optional; they are necessary for coexistence.
\end{corollary}

\textbf{Interpretation:} Unconstrained coexistence is impossible in any
shared finite-resource environment. An agent can achieve its
unconstrained optimum only if no other agent competes for any of the
same resources.

\section{The Multi-Agent Coupled System: Generalized Nash
Equilibrium}\label{the-multi-agent-coupled-system-generalized-nash-equilibrium}

\subsection{Formulation}\label{formulation-1}

When all \(N\) agents optimize simultaneously, each taking the others'
strategies into account, the system is a \textbf{Generalized Nash
Equilibrium Problem (GNEP)}. Each agent \(i\) solves:

\[\max_{\mathbf{x}_i \geq \mathbf{0}} \; U_i(\mathbf{x}_i) \qquad \text{s.t.} \quad \sum_{l=1}^{N} x_{lj} \leq R_j \;\; \forall j, \quad g_{Bk}(\mathbf{x}_i) \leq 0 \;\; \forall B \neq i, \; k\]

The feasible set for each agent depends on the other agents' strategies
(through the shared resource constraints). This is what makes the
problem a \emph{generalized} Nash problem, not a standard one.

\subsection{Solution Concept: Variational
Equilibrium}\label{solution-concept-variational-equilibrium}

A \textbf{Variational Equilibrium} (VE), also called a Normalized Nash
Equilibrium \citep{Rosen1965}, is a strategy profile
\((\mathbf{x}_1^*, \ldots, \mathbf{x}_N^*)\) such that:

\begin{enumerate}
\def\labelenumi{\arabic{enumi}.}
\tightlist
\item
  Each agent \(i\)'s strategy \(\mathbf{x}_i^*\) is optimal given the
  others' strategies, and
\item
  All agents share the \textbf{same} Lagrange multipliers
  \(\boldsymbol{\lambda}^* = (\lambda_1^*, \ldots, \lambda_n^*)\) for
  the shared resource constraints.
\end{enumerate}

\begin{theorem}[Existence and Structure of the Variational Equilibrium]
\label{thm-variational-equilibrium}

Under Assumption~\ref{asm-regularity}, the GNEP possesses a Variational Equilibrium $(\mathbf{x}_1^*, \ldots, \mathbf{x}_N^*, \boldsymbol{\lambda}^*)$. At this equilibrium:

(a) The shared shadow prices $\lambda_j^*$ are identical for all agents: every agent faces the same "price" for each scarce resource.

(b) The stationarity condition for each agent $i$ reads:

$$\frac{\partial U_i}{\partial x_{ij}}\bigg|_{\mathbf{x}_i^*} = \lambda_j^* + \sum_{B \neq i, k} \mu_{Bk}^{(i)} \frac{\partial g_{Bk}}{\partial x_{ij}}\bigg|_{\mathbf{x}_i^*} \qquad \forall j$$

(c) No agent can improve its utility by unilaterally deviating from $\mathbf{x}_i^*$, given the constraints.
\end{theorem}

\begin{proof}
The shared resource constraints $0 \leq x_{ij} \leq R_j$ give each agent a compact, convex feasible set, each $U_i$ is concave by Assumption~\ref{asm-regularity}, and the joint constraint set is convex. These are the hypotheses of \citep[Theorem 3]{Rosen1965}, which guarantees that a Normalized Nash Equilibrium exists for every weight vector $r > 0$. The shared-multiplier property is the defining feature of the VE solution concept.
\end{proof}

\subsection{Interpretation: Equality Before the
Constraints}\label{interpretation-equality-before-the-constraints}

\begin{corollary}[Symmetry of Scarcity]
\label{cor-symmetry-scarcity}

At the variational equilibrium, all agents experience the same marginal cost $\lambda_j^*$ per unit of each shared resource. No agent is privileged in the structure of the shared constraints.
\end{corollary}

The scarcity constraint imposes a uniform shadow-price structure on all
agents. Individual differences arise only through differences in utility
functions \(U_i\) (different agents value resources differently) and
through the specific boundary constraints \(g_{Bk}\) (different agents
have different protected domains).

\subsection{Welfare Optimality of the Variational
Equilibrium}\label{welfare-optimality-of-the-variational-equilibrium}

\begin{theorem}[Welfare Optimality]
\label{thm-welfare-optimality}

Under Assumption~\ref{asm-regularity}, the variational equilibrium $(\mathbf{x}_1^*, \ldots, \mathbf{x}_N^*)$ is the unique maximizer of total social welfare over the feasible set:

$$\sum_{i=1}^{N} U_i(\mathbf{x}_i^*) > \sum_{i=1}^{N} U_i(\mathbf{x}_i) \qquad \forall \; (\mathbf{x}_1, \ldots, \mathbf{x}_N) \in \mathcal{F} \setminus \{\mathbf{x}^*\}$$

That is, the VE uniquely maximizes $W(\mathbf{x}) = \sum_{i=1}^{N} U_i(\mathbf{x}_i)$ subject to the shared resource and boundary constraints.
\end{theorem}

\begin{proof}
The argument is by direct KKT equivalence. Because each $U_i$ depends only on agent $i$'s own action vector $\mathbf{x}_i$ (coupling between agents enters solely through the shared resource and boundary constraints, not through the utilities), the gradient of the social welfare objective decomposes agent-wise:

$$\frac{\partial W}{\partial x_{ij}} = \frac{\partial U_i}{\partial x_{ij}} \qquad \forall i, j.$$

Writing the KKT conditions for the welfare problem $\max_{\mathbf{x} \in \mathcal{F}} \sum_i U_i(\mathbf{x}_i)$ with multipliers $\lambda_j$ on the shared resource constraints and $\mu_{Bk}^{(i)}$ on the boundary constraints that restrict agent $i$'s action yields, for every $i, j$:

$$\frac{\partial U_i}{\partial x_{ij}} = \lambda_j + \sum_{B \neq i, k} \mu_{Bk}^{(i)} \frac{\partial g_{Bk}}{\partial x_{ij}},$$

together with primal/dual feasibility and complementary slackness. These are \emph{identical} to the VE stationarity conditions of Theorem Theorem~\ref{thm-variational-equilibrium}(b), with the shared resource multipliers $\lambda_j$ playing the role of the common shadow prices $\lambda_j^*$ that defines the variational (normalized) equilibrium.

Under Assumption~\ref{asm-regularity} and Slater's condition (established in the KKT discussion above), these conditions are both necessary and sufficient for optimality of the welfare problem \citep[§5.5.3]{Boyd2004}. Hence the VE allocation satisfies the welfare-maximization KKT system, and conversely any welfare maximizer is a VE.

Uniqueness follows from strict concavity: each $U_i$ is strictly concave by Assumption~\ref{asm-regularity}(b), so $W = \sum_i U_i$ is strictly concave on the convex, compact feasible set $\mathcal{F}$, admitting at most one maximizer. Existence is guaranteed by continuity of $W$ on $\mathcal{F}$. Therefore the VE is the unique global maximizer of total social welfare over $\mathcal{F}$. (This uniform-weight coincidence between the normalized Nash equilibrium and the social welfare optimum is the separable-utility case of \citep[§5]{Rosen1965}.)
\end{proof}

\begin{corollary}[Welfare Dominance]
\label{cor-welfare-dominance}

The variational equilibrium is welfare-superior to every other feasible allocation, including equal-split, greedy (unconstrained-optimal for a single agent), and any arbitrary feasible partition of resources.
\end{corollary}

\subsection{The ``Mutual Constraint''
Structure}\label{the-mutual-constraint-structure}

At the variational equilibrium, the full system of KKT conditions across
all agents forms a coupled system:

\[\frac{\partial U_i}{\partial x_{ij}} = \lambda_j^* + \sum_{B \neq i, k} \mu_{Bk}^{(i)} \frac{\partial g_{Bk}}{\partial x_{ij}} \qquad \forall i, j\]

\[\sum_{i=1}^{N} x_{ij}^* \leq R_j, \quad \lambda_j^* \geq 0, \quad \lambda_j^*\left(\sum_i x_{ij}^* - R_j\right) = 0 \qquad \forall j\]

\[g_{Bk}(\mathbf{x}_i^*) \leq 0, \quad \mu_{Bk}^{(i)} \geq 0, \quad \mu_{Bk}^{(i)} g_{Bk}(\mathbf{x}_i^*) = 0 \qquad \forall i, B \neq i, k\]

This coupled system is a self-consistent set of mutual constraints that
no agent can improve upon unilaterally, a fixed point of the multi-agent
optimization. The system of mutual boundary constraints is not imposed
exogenously; it emerges as the necessary structure of any stable
multi-agent coexistence under scarcity.

\section{Worked Example: Two Agents, One Divisible
Resource}\label{worked-example-two-agents-one-divisible-resource}

\subsection{Setup}\label{setup}

Let \(N = 2\), \(n = 1\) (agents \(A\) and \(B\), one resource of total
quantity \(R\)). The quadratic utility functions are:

\[U_A(x_A) = \alpha_A x_A - \frac{\beta_A}{2} x_A^2, \qquad U_B(x_B) = \alpha_B x_B - \frac{\beta_B}{2} x_B^2\]

where \(x_B = R - x_A\) (by resource conservation, taking the constraint
as binding for simplicity, since this is the interesting case).

\subsection{Unconstrained Optima}\label{unconstrained-optima}

\[x_A^{\circ} = \frac{\alpha_A}{\beta_A}, \qquad x_B^{\circ} = \frac{\alpha_B}{\beta_B}\]

\textbf{Collision condition:} \(x_A^{\circ} + x_B^{\circ} > R\), i.e.,
\(\frac{\alpha_A}{\beta_A} + \frac{\alpha_B}{\beta_B} > R\).

\subsection{\texorpdfstring{With Boundary Constraints: Agent \(B\) Has a
Guaranteed Minimum
\(\bar{x}_B\)}{With Boundary Constraints: Agent B Has a Guaranteed Minimum \textbackslash bar\{x\}\_B}}\label{with-boundary-constraints-agent-b-has-a-guaranteed-minimum-barx_b}

Agent \(A\) maximizes \(U_A(x_A)\) subject to
\(x_A \leq R - \bar{x}_B\).

\textbf{Lagrangian:}

\[\mathcal{L}(x_A, \mu) = \alpha_A x_A - \frac{\beta_A}{2} x_A^2 - \mu\left(x_A - R + \bar{x}_B\right)\]

\textbf{KKT conditions:}

\[\alpha_A - \beta_A x_A^* - \mu^* = 0\]
\[\mu^* \geq 0, \quad x_A^* \leq R - \bar{x}_B, \quad \mu^*(x_A^* - R + \bar{x}_B) = 0\]

\textbf{Solution:}

If \(\frac{\alpha_A}{\beta_A} \leq R - \bar{x}_B\) (constraint not
binding): \[x_A^* = \frac{\alpha_A}{\beta_A}, \qquad \mu^* = 0\]

If \(\frac{\alpha_A}{\beta_A} > R - \bar{x}_B\) (constraint binding):
\[x_A^* = R - \bar{x}_B, \qquad \mu^* = \alpha_A - \beta_A(R - \bar{x}_B)\]

\subsection{Numerical Illustration}\label{numerical-illustration}

Let \(\alpha_A = 10\), \(\beta_A = 1\), \(\alpha_B = 8\),
\(\beta_B = 1\), \(R = 12\), \(\bar{x}_B = 5\).

\begin{itemize}
\tightlist
\item
  Unconstrained optima: \(x_A^{\circ} = 10\), \(x_B^{\circ} = 8\).
  Collision: \(10 + 8 = 18 > 12\).
\item
  With \(B\)'s boundary constraint: \(A\) can take at most
  \(R - \bar{x}_B = 7\).
\item
  Since \(x_A^{\circ} = 10 > 7\), the constraint binds: \(x_A^* = 7\),
  \(\mu^* = 10 - 1 \cdot 7 = 3\).
\item
  \textbf{The cost of \(B\)'s boundary constraint to \(A\):} Agent \(A\)
  loses \(\mu^* = 3\) units of marginal energy due to \(B\)'s guaranteed
  share.
\item
  For these parameters the choice \(\bar{x}_B = 5\) recovers the
  variational equilibrium (\(\lambda^* = (10+8-12)/(1+1) = 3\), giving
  \(x_A^* = 7\), \(x_B^* = 5\)). The welfare comparisons that follow
  therefore illustrate Theorem~\ref{thm-welfare-optimality} rather than
  a coincidence of parameter choice.
\item
  Agent \(A\)'s constrained utility:
  \(U_A(7) = 10(7) - \frac{1}{2}(49) = 70 - 24.5 = 45.5\).
\item
  Agent \(A\)'s unconstrained utility:
  \(U_A(10) = 10(10) - \frac{1}{2}(100) = 100 - 50 = 50\).
\item
  \textbf{Total cost to \(A\):} \(50 - 45.5 = 4.5\) energy units.
\end{itemize}

\subsection{Social Welfare Comparison}\label{social-welfare-comparison}

\textbf{With boundary constraints (\(\bar{x}_B = 5\)):}

\begin{itemize}
\tightlist
\item
  \(U_A(7) = 45.5\), \(U_B(5) = 8(5) - \frac{1}{2}(25) = 27.5\).
\item
  \textbf{Total system utility:} \(45.5 + 27.5 = 73.0\).
\end{itemize}

\textbf{Without boundary constraints (collision → conflict, modeled as
equal split \(x_A = x_B = 6\)):}

\begin{itemize}
\tightlist
\item
  \(U_A(6) = 10(6) - \frac{1}{2}(36) = 42\),
  \(U_B(6) = 8(6) - \frac{1}{2}(36) = 30\).
\item
  \textbf{Total system utility:} \(42 + 30 = 72.0\).
\end{itemize}

\textbf{Without boundary constraints (A takes all, unconstrained):}

\begin{itemize}
\tightlist
\item
  \(U_A(10) = 50\), \(U_B(2) = 8(2) - \frac{1}{2}(4) = 14\).
\item
  \textbf{Total system utility:} \(50 + 14 = 64.0\).
\end{itemize}

The boundary-constrained allocation achieves the highest total system
utility. The unconstrained ``greedy'' solution is the worst for the
system.

\begin{quote}
\textbf{Remark (Why the gap is small here).} The welfare difference
between the VE allocation (73.0) and the naive equal split (72.0) is
only \textasciitilde1.4\%. This is a \emph{mathematical artifact of
near-symmetric agents}: when \(\beta_A = \beta_B\), the VE allocation
\(x_i^* = (\alpha_i - \lambda^*)/\beta\) differs from equal split only
through the \(\alpha_i\) differences. With \(\alpha_A = 10\) and
\(\alpha_B = 8\) (a 20\% gap), the heterogeneity is small and equal
split is nearly optimal. In general, the welfare gain from optimally
differentiated allocation scales with agent heterogeneity. The following
asymmetric example demonstrates this.
\end{quote}

\subsection{Asymmetric Example: Heterogeneous Metabolic
Efficiency}\label{asymmetric-example-heterogeneous-metabolic-efficiency}

In biological systems, agents rarely have identical metabolic profiles.
Consider two agents with the \textbf{same marginal yield} (\(\alpha\))
but \textbf{different diminishing-returns rates} (\(\beta\)): for
example, Agent \(A\) is an efficient large-scale processor (forests,
apex predators, industrial economies) while Agent \(B\) saturates
quickly (specialized microorganisms, artisanal producers).

\textbf{Parameters:} \(\alpha_A = 10\), \(\beta_A = 0.5\) (slow
diminishing returns); \(\alpha_B = 10\), \(\beta_B = 2\) (fast
diminishing returns); \(R = 12\).

\textbf{Unconstrained optima:} \(x_A^{\circ} = 20\),
\(x_B^{\circ} = 5\). Collision: \(25 > 12\).

\textbf{Variational Equilibrium:}

\[\lambda^* = \frac{x_A^{\circ} + x_B^{\circ} - R}{1/\beta_A + 1/\beta_B} = \frac{20 + 5 - 12}{2 + 0.5} = \frac{13}{2.5} = 5.2\]

\[x_A^* = \frac{10 - 5.2}{0.5} = 9.6, \qquad x_B^* = \frac{10 - 5.2}{2} = 2.4\]

The VE allocates \textbf{four times more} to the efficient processor,
not because \(A\) is ``more important,'' but because \(A\) extracts more
marginal value from each additional unit.

\textbf{Social welfare comparison:}

{\def\LTcaptype{none} % do not increment counter
\begin{longtable}[]{@{}
  >{\raggedright\arraybackslash}p{(\linewidth - 10\tabcolsep) * \real{0.1667}}
  >{\raggedright\arraybackslash}p{(\linewidth - 10\tabcolsep) * \real{0.1667}}
  >{\raggedright\arraybackslash}p{(\linewidth - 10\tabcolsep) * \real{0.1667}}
  >{\raggedright\arraybackslash}p{(\linewidth - 10\tabcolsep) * \real{0.1667}}
  >{\raggedright\arraybackslash}p{(\linewidth - 10\tabcolsep) * \real{0.1667}}
  >{\raggedright\arraybackslash}p{(\linewidth - 10\tabcolsep) * \real{0.1667}}@{}}
\toprule\noalign{}
\begin{minipage}[b]{\linewidth}\raggedright
Allocation
\end{minipage} & \begin{minipage}[b]{\linewidth}\raggedright
\(x_A\)
\end{minipage} & \begin{minipage}[b]{\linewidth}\raggedright
\(x_B\)
\end{minipage} & \begin{minipage}[b]{\linewidth}\raggedright
\(U_A\)
\end{minipage} & \begin{minipage}[b]{\linewidth}\raggedright
\(U_B\)
\end{minipage} & \begin{minipage}[b]{\linewidth}\raggedright
\textbf{Total SW}
\end{minipage} \\
\midrule\noalign{}
\endhead
\bottomrule\noalign{}
\endlastfoot
\textbf{VE (optimal boundary constraints)} & 9.6 & 2.4 & 72.96 & 18.24 &
\textbf{91.20} \\
\textbf{Greedy (\(A\) takes all)} & 12 & 0 & 84.00 & 0.00 & 84.00 \\
\textbf{Equal split} & 6 & 6 & 51.00 & 24.00 & 75.00 \\
\end{longtable}
}

The VE now outperforms equal split by \textbf{21.6\%} and the greedy
allocation by \textbf{8.6\%}. The equal-split regime is particularly
wasteful: it forces 6 units onto Agent \(B\), whose fast-saturating
metabolism means each unit beyond \textasciitilde2.4 costs more to
process than it yields. Biologically, this is analogous to force-feeding
an organism past its metabolic optimum: raw resource input does not
equate to usable energy output.

\textbf{Key insight:} The framework does not merely prove that
\emph{some} constraint system beats anarchy. It proves that the
\emph{optimal} constraint system (the VE) allocates resources according
to each agent's capacity to convert them into useful work, and this
allocation can dramatically outperform equal-split allocation precisely
when agents are diverse.

\subsection{Symmetric Example: Same-Species Competition and the Scarcity
Gradient}\label{symmetric-example-same-species-competition-and-the-scarcity-gradient}

A ubiquitous biological scenario is intraspecific competition: agents of
the \textbf{same species} sharing the \textbf{same environment}
(\(\alpha_A = \alpha_B = \alpha\), \(\beta_A = \beta_B = \beta\)). Here
the framework shows a different but equally important result.

\textbf{Parameters:} \(\alpha = 10\), \(\beta = 1\); resource \(R\)
varies.

\textbf{Unconstrained optima:} \(x_A^{\circ} = x_B^{\circ} = 10\).
\textbf{Collision} occurs whenever \(R < 20\).

\textbf{VE for symmetric agents:} By symmetry of the KKT system, the
variational equilibrium assigns each agent exactly half:

\[x_A^* = x_B^* = \frac{R}{2}, \qquad \lambda^* = \alpha - \beta \cdot \frac{R}{2}\]

\textbf{Equal allocation is not assumed; it is derived.} The equal split
emerges as a theorem from the symmetry of identical utility functions
under shared scarcity constraints. No exogenous distributional rule is
required.

\textbf{The cost of defection:} If one agent defects (seizes its
unconstrained optimum \(x_A = 10\), leaving the other with \(R - 10\)),
the total system welfare drops. The magnitude of this loss scales with
scarcity:

{\def\LTcaptype{none} % do not increment counter
\begin{longtable}[]{@{}
  >{\raggedright\arraybackslash}p{(\linewidth - 14\tabcolsep) * \real{0.1250}}
  >{\raggedright\arraybackslash}p{(\linewidth - 14\tabcolsep) * \real{0.1250}}
  >{\raggedright\arraybackslash}p{(\linewidth - 14\tabcolsep) * \real{0.1250}}
  >{\raggedright\arraybackslash}p{(\linewidth - 14\tabcolsep) * \real{0.1250}}
  >{\raggedright\arraybackslash}p{(\linewidth - 14\tabcolsep) * \real{0.1250}}
  >{\raggedright\arraybackslash}p{(\linewidth - 14\tabcolsep) * \real{0.1250}}
  >{\raggedright\arraybackslash}p{(\linewidth - 14\tabcolsep) * \real{0.1250}}
  >{\raggedright\arraybackslash}p{(\linewidth - 14\tabcolsep) * \real{0.1250}}@{}}
\toprule\noalign{}
\begin{minipage}[b]{\linewidth}\raggedright
\(R\)
\end{minipage} & \begin{minipage}[b]{\linewidth}\raggedright
Scarcity
\end{minipage} & \begin{minipage}[b]{\linewidth}\raggedright
VE: \(x_i^*\)
\end{minipage} & \begin{minipage}[b]{\linewidth}\raggedright
VE: \(U_i\)
\end{minipage} & \begin{minipage}[b]{\linewidth}\raggedright
VE: SW
\end{minipage} & \begin{minipage}[b]{\linewidth}\raggedright
Defect: \((x_A, x_B)\)
\end{minipage} & \begin{minipage}[b]{\linewidth}\raggedright
Defect: SW
\end{minipage} & \begin{minipage}[b]{\linewidth}\raggedright
\textbf{VE Advantage}
\end{minipage} \\
\midrule\noalign{}
\endhead
\bottomrule\noalign{}
\endlastfoot
20 & None & 10.0 & 50.0 & 100.0 & (10, 10) & 100.0 & 0\% \\
16 & Mild & 8.0 & 48.0 & 96.0 & (10, 6) & 92.0 & \textbf{4.3\%} \\
14 & Moderate & 7.0 & 45.5 & 91.0 & (10, 4) & 82.0 & \textbf{11.0\%} \\
12 & Scarce & 6.0 & 42.0 & 84.0 & (10, 2) & 68.0 & \textbf{23.5\%} \\
10 & Severe & 5.0 & 37.5 & 75.0 & (10, 0) & 50.0 & \textbf{50.0\%} \\
\end{longtable}
}

Three results emerge from this table:

\begin{enumerate}
\def\labelenumi{\arabic{enumi}.}
\item
  \textbf{Under abundance (\(R \geq 20\)), boundary constraints have
  zero cost}: both agents achieve their unconstrained optima, all
  multipliers are zero, and the constraint system is slack. Boundary
  constraints exist in principle but impose no restriction. This is the
  formal counterpart of the observation that resource conflicts rarely
  arise when resources are plentiful.
\item
  \textbf{As scarcity increases, the cost of defection grows
  super-linearly.} Moving from mild to severe scarcity, the VE advantage
  jumps from 4\% to 50\%. This is because the defector's gain (moving
  from \(R/2\) to 10) is bounded by the concavity of \(U\), while the
  victim's loss (moving from \(R/2\) to \(R - 10\)) accelerates as it
  approaches zero, the region where each lost unit is maximally
  valuable.
\item
  \textbf{Under severe scarcity (\(R = 10\)), defection eliminates the
  other agent entirely} (\(x_B = 0 \implies U_B = 0\)). The defector
  captures \(U_A = 50\) (a gain of 12.5 over the VE), but agent \(B\)
  loses its entire share of 37.5, yielding a net welfare destruction of
  25 units. This is the signature of Thermodynamic Friction
  \citep{TC-IV}: unrestricted competition in scarce environments does
  not redistribute value, it \emph{dissipates} it.
\end{enumerate}

\begin{corollary}[Boundary Constraint Value Scales with Scarcity]
\label{cor-constraint-scarcity-scaling}

The value of a boundary constraint system is monotonically increasing in scarcity. Under abundance, boundary constraints are costless and unnecessary. Under scarcity, boundary constraints become the dominant determinant of system welfare.
\end{corollary}

The scarcity-monotonicity result provides a mathematical basis for the
prediction that multi-agent systems will intensify coordination under
resource pressure: the welfare cost of failed coordination grows with
scarcity.

\section{The Transition to Conflict: Removing the Boundary
Constraints}\label{the-transition-to-conflict-removing-the-boundary-constraints}

\subsection{Statement}\label{statement}

\begin{proposition}[Unconstrained Regime Implies Conflict]
\label{prop-unconstrained-conflict}

If all boundary constraints are removed ($g_{Bk}$ dropped for all $B, k$), and the system is in collision ($\sum_i x_{ij}^{\circ} > R_j$ for some $j$), then:

(a) No stable allocation exists without an enforcement mechanism.

(b) Each agent's attempt to realize its unconstrained optimum generates a negative externality on other agents.

(c) The resulting contested resource must be resolved by direct energy expenditure (conflict), creating the interaction cost function quantified in the thermodynamic friction derivation.
\end{proposition}

\begin{proof}
\emph{Proof of (a).} Without constraints, each agent $i$ attempts $\mathbf{x}_i^{\circ}$. The aggregate $\sum_i \mathbf{x}_i^{\circ} > \mathbf{R}$ is physically infeasible. Since no constraint directs agents to reduce consumption, and each agent's unilateral best response is $\mathbf{x}_i^{\circ}$ regardless of the others' (by strict concavity, the optimal response function without constraints is constant), there is no self-correcting mechanism. The unconstrained Nash equilibrium is infeasible: it requires $\sum_i \mathbf{x}_i^{\circ} > \mathbf{R}$. Any feasible allocation requires at least one agent to accept less than its optimum, which it has no incentive to do without a constraint. $\square$

\emph{Proof of (b).} When $A$ claims $x_{Aj}^{\circ}$, the remaining resource for all other agents is $R_j - x_{Aj}^{\circ}$. In the collision regime, $R_j - x_{Aj}^{\circ} < \sum_{i \neq A} x_{ij}^{\circ}$, meaning $A$'s strategy forces at least one other agent below its optimum. This is the definition of a negative externality.
\end{proof}

\textbf{Remark.} Part (c) connects directly to the thermodynamic
friction derivation \citep{TC-IV}. When the constraint is absent, the
physical reality of resource exclusion (two agents cannot consume the
same molecule) is enforced not by coordination but by direct energy
expenditure: adversarial contest over the disputed resource. The
constraint \(g_{Bk}(\mathbf{x}_A) \leq 0\) replaces costly physical
enforcement with costless (zero-friction) boundary recognition.

\section{Summary of Results}\label{summary-of-results}

{\def\LTcaptype{none} % do not increment counter
\begin{longtable}[]{@{}
  >{\raggedright\arraybackslash}p{(\linewidth - 6\tabcolsep) * \real{0.2500}}
  >{\raggedright\arraybackslash}p{(\linewidth - 6\tabcolsep) * \real{0.2500}}
  >{\raggedright\arraybackslash}p{(\linewidth - 6\tabcolsep) * \real{0.2500}}
  >{\raggedright\arraybackslash}p{(\linewidth - 6\tabcolsep) * \real{0.2500}}@{}}
\toprule\noalign{}
\begin{minipage}[b]{\linewidth}\raggedright
\#
\end{minipage} & \begin{minipage}[b]{\linewidth}\raggedright
Result
\end{minipage} & \begin{minipage}[b]{\linewidth}\raggedright
Statement
\end{minipage} & \begin{minipage}[b]{\linewidth}\raggedright
Significance
\end{minipage} \\
\midrule\noalign{}
\endhead
\bottomrule\noalign{}
\endlastfoot
\textbf{T1} & Shadow Price Theorem &
\(\partial U_A^* / \partial c_{Bk} = \mu_{Bk}^* \geq 0\) & Every
boundary constraint has a quantifiable energetic cost to the constrained
agent \\
\textbf{C1.1} & Active vs.~Inactive boundary constraints & Binding
constraints have \(\mu > 0\); slack constraints have \(\mu = 0\) &
Boundary constraints cost nothing when not contested; cost emerges only
at points of conflict \\
\textbf{C1.2} & Coupled Cost of Active Constraints & Every active
boundary constraint has \(\mu_{Bk}^* > 0\), simultaneously protecting
\(B\) and restricting \(A\) & The same multiplier measures protection
and restriction \\
\textbf{L1} & Inevitability of Collision & Collision guaranteed for
\(N \geq N_j^*\) agents sharing finite resources & Scarcity makes
constraint necessity quantitative, not merely qualitative \\
\textbf{T2} & Impossibility of Infinite Freedom & No feasible profile
achieves all agents' unconstrained optima simultaneously & Boundary
constraints are necessary, not optional \\
\textbf{C2.1} & Necessity of Active Constraints & At least one
\(\mu_{Bk}^* > 0\) at any feasible solution & Some agent is always
constrained \\
\textbf{T3} & Variational Equilibrium & GNEP has a solution with shared
shadow prices & Uniform shadow-price structure emerges from the shared
constraint structure \\
\textbf{C3.1} & Symmetry of Scarcity & All agents face the same
\(\lambda_j^*\) per shared resource & Uniform cost structure; individual
differences arise only through \(U_i\) and \(g_{Bk}\) \\
\textbf{T4} & Welfare Optimality & VE uniquely maximizes total social
welfare over the feasible set & The VE is the best feasible allocation,
not merely a stable one \\
\textbf{C4.1} & Welfare Dominance & VE is welfare-superior to every
other feasible allocation & Greedy, equal-split, and all other
alternatives are strictly dominated \\
\textbf{C4.2} & Boundary Constraint Value Scales with Scarcity & Value
of boundary constraints monotonically increasing in scarcity &
Coordination intensifies during crises as a consequence of the model \\
\textbf{P1} & Unconstrained → Conflict & Removing constraints in
collision regime produces no stable allocation & The unconstrained
regime is unstable; links to thermodynamic friction \\
\end{longtable}
}

\section{Notation Index}\label{notation-index}

{\def\LTcaptype{none} % do not increment counter
\begin{longtable}[]{@{}
  >{\raggedright\arraybackslash}p{(\linewidth - 2\tabcolsep) * \real{0.5000}}
  >{\raggedright\arraybackslash}p{(\linewidth - 2\tabcolsep) * \real{0.5000}}@{}}
\toprule\noalign{}
\begin{minipage}[b]{\linewidth}\raggedright
Symbol
\end{minipage} & \begin{minipage}[b]{\linewidth}\raggedright
Meaning
\end{minipage} \\
\midrule\noalign{}
\endhead
\bottomrule\noalign{}
\endlastfoot
\(N\) & Number of agents \\
\(n\) & Number of resource dimensions \\
\(\mathbf{x}_i \in \mathbb{R}_{\geq 0}^n\) & Agent \(i\)'s strategy
(resource-allocation) vector \\
\(R_j\) & Total available quantity of resource \(j\) \\
\(U_i(\mathbf{x}_i)\) & Agent \(i\)'s utility (boundary-maintenance
objective) \\
\(\alpha_{ij}, \beta_{ij}\) & Marginal yield and diminishing-returns
parameters (quadratic model) \\
\(\mathbf{x}_i^{\circ}\) & Agent \(i\)'s unconstrained optimum \\
\(g_{Bk}(\mathbf{x}_A)\) & Constraint function encoding agent \(B\)'s
\(k\)-th boundary constraint \\
\(\lambda_j\) & Lagrange multiplier for resource \(j\) scarcity \\
\(\mu_{Bk}\) & Lagrange multiplier (shadow price) for agent \(B\)'s
\(k\)-th boundary constraint \\
\(\hat{x}_{-Aj}\) & Aggregate consumption of resource \(j\) by all
agents other than \(A\) \\
\(\mathcal{L}_A\) & Lagrangian for agent \(A\)'s constrained
optimization \\
\(\mathcal{F}\) & Feasible strategy set under resource and boundary
constraints \\
\end{longtable}
}

\begin{quote}
\textbf{Symbol disambiguation.} The symbol \(R_j\) denotes total
resource quantity. In subsequent papers in this series, \(r\) denotes
coupling distance \citep{TC-VI}, contest decisiveness \citep{TC-IV}, and
time-preference rate \citep{TC-V}. Context and subscripts distinguish
these uses.
\end{quote}

\section{Discussion}\label{discussion}

\subsection{Physical Interpretation}\label{physical-interpretation-1}

The results form a complete logical chain. Scarcity creates collision
(Lemma~\ref{lem-inevitability-collision}); collision requires boundary
constraints (Theorem~\ref{thm-impossibility-infinite-freedom}); the
constrained system admits a variational equilibrium with uniform shadow
prices (Theorem~\ref{thm-variational-equilibrium}); and constraint value
scales monotonically with scarcity
(Corollary~\ref{cor-constraint-scarcity-scaling}).

Each result is denominated in energy units and computationally verified.
The shadow prices of Theorem~\ref{thm-shadow-price} provide an
auditable, real-time signal of constraint tension, a computable measure
of how costly a given constraint is to the constrained agent.

The central contribution is the demonstration that boundary constraints
in multi-agent dissipative systems are not design choices or social
conventions but necessities: the GNEP has no feasible solution without
them.

\subsection{Testable Predictions}\label{testable-predictions}

The model generates falsifiable predictions with explicit failure
conditions:

{\def\LTcaptype{none} % do not increment counter
\begin{longtable}[]{@{}
  >{\raggedright\arraybackslash}p{(\linewidth - 4\tabcolsep) * \real{0.3333}}
  >{\raggedright\arraybackslash}p{(\linewidth - 4\tabcolsep) * \real{0.3333}}
  >{\raggedright\arraybackslash}p{(\linewidth - 4\tabcolsep) * \real{0.3333}}@{}}
\toprule\noalign{}
\begin{minipage}[b]{\linewidth}\raggedright
Prediction
\end{minipage} & \begin{minipage}[b]{\linewidth}\raggedright
Source
\end{minipage} & \begin{minipage}[b]{\linewidth}\raggedright
Falsified If\ldots{}
\end{minipage} \\
\midrule\noalign{}
\endhead
\bottomrule\noalign{}
\endlastfoot
Uniform shadow prices across agents for each shared resource at the VE &
Theorem~\ref{thm-variational-equilibrium} & Resource markets with
physical delivery constraints consistently exhibit agent-specific
(non-uniform) shadow prices under conditions satisfying
Assumption~\ref{asm-regularity} \\
Welfare gain from constraints is monotonically increasing in scarcity &
Corollary~\ref{cor-constraint-scarcity-scaling} & Cooperation rates or
welfare measures decrease as resource scarcity increases, controlling
for other variables \\
Unconstrained regime is strictly dominated for all scarcity levels
\(R < \sum_i x_i^{\circ}\) &
Theorem~\ref{thm-impossibility-infinite-freedom}, worked examples
(§\ref{worked-example-two-agents-one-divisible-resource}) & Experimental
resource-sharing games with energy-denominated payoffs find that
unconstrained allocation outperforms the constrained VE allocation \\
\end{longtable}
}

The strongest falsification of the model's core result would be a class
of physically realistic multi-agent systems, with concave utility
functions and finite shared resources satisfying
Assumption~\ref{asm-regularity}, in which a feasible allocation exists
that is simultaneously unconstrained-optimal for all agents. The model's
architecture makes this structurally impossible under resource
collision, but a counterexample would require revision of
Theorem~\ref{thm-impossibility-infinite-freedom}.

\subsection{Limitations}\label{limitations}

The model inherits standard idealizations. The utility functions satisfy
regularity conditions (Assumption~\ref{asm-regularity}) that may not
hold for all physical systems; in particular, strict concavity and
coercivity exclude agents with increasing returns to scale or threshold
effects. The two-agent worked examples illustrate the general
\(N\)-agent theorems but do not capture heterogeneous or large-\(N\)
populations. The static resource endowment assumption excludes
renewable, seasonal, or stochastic resource environments. The
restriction to negative boundary constraints
(\(g_{Bk}(\mathbf{0}) < 0\)) means the model addresses non-interference
but not provision obligations (see §\ref{future-work}).

The constraint necessity result is a feasibility argument: it shows that
no feasible allocation exists without constraints, but does not specify
the mechanism by which constraints are established or enforced.
\emph{Thermodynamic Friction} \citep{TC-IV} quantifies the energy cost
of constraint violation, and \emph{Cooperative Equilibrium} \citep{TC-V}
derives game-theoretic enforcement via repeated interaction.

\textbf{Scope.} The model applies to any system of agents satisfying
Assumption~\ref{asm-regularity} that share finite resources, without
restriction on species, substrate, or scale. But the results address the
\emph{structural rules} of multi-agent allocation: constraint necessity,
shadow pricing, and equilibrium existence. Questions of distributive
justice beyond Pareto efficiency, the political process by which
specific constraint levels are negotiated, and the enforcement
architecture that maintains constraints in practice lie outside the
optimization problem this paper defines.

\subsection{Applications}\label{applications}

The model applies wherever multiple self-maintaining systems share
finite resources:

\begin{enumerate}
\def\labelenumi{\arabic{enumi}.}
\item
  \textbf{Ecological carrying capacity.} Species competing for
  overlapping resource niches face the collision condition of
  Lemma~\ref{lem-inevitability-collision}. The boundary constraints
  correspond to territorial boundaries, niche partitioning, and
  behavioral dominance hierarchies, mechanisms that ecologists observe
  empirically. The model proves that some constraint system is necessary
  for stable coexistence under resource collision; these ecological
  mechanisms are empirical instances of such constraints. The
  scarcity-scaling corollary predicts that niche-partitioning intensity
  should increase with resource scarcity.
\item
  \textbf{Common-pool resource management.} \citet{Ostrom1990}
  documented institutional constraint systems (harvest limits, access
  rules, monitoring) that communities develop to manage shared
  fisheries, irrigation systems, and forests. These are instances of
  mutually enforced boundary constraints. The scarcity-scaling corollary
  (Corollary~\ref{cor-constraint-scarcity-scaling}) predicts that such
  institutions intensify as the resource base declines.
\item
  \textbf{Economic markets.} The variational equilibrium's uniform
  shadow prices are a mathematical analog of competitive market prices.
  Within the model, the shared Lagrange multipliers play the role
  conventionally assigned to market-clearing prices, providing one
  possible microfoundation that does not invoke an auctioneer mechanism.
  The worked examples
  (§\ref{worked-example-two-agents-one-divisible-resource}) show that
  the VE allocation outperforms both greedy and equal-split regimes,
  with the welfare advantage increasing with agent heterogeneity.
\end{enumerate}

\subsection{Future Work}\label{future-work}

\textbf{Positive boundary constraints and endogenous production.} The
model is restricted to non-interference conditions
(\(g_{Bk}(\mathbf{0}) < 0\)). Provision constraints, where
\(g_{Bk}(\mathbf{0}) > 0\) so that inaction itself violates the
constraint, arise as stability conditions on an already-established
variational equilibrium. Formally, this requires replacing
\(g_{Bk}(\mathbf{0}) < 0\) with direct Slater's condition (non-empty
feasible-set interior) and deriving when equilibrium maintenance
requires active resource transfer toward vulnerable agents. The
historical pattern in which non-interference constraints tend to precede
provision constraints in legal and institutional development is
consistent with the ordering the model's constraint hierarchy predicts.

Extending the model to a production network in which agents' outputs
feed other agents' inputs would convert positive boundary constraints
into a derived result: whenever agent \(B\)'s output is essential input
for agent \(A\), the provision threshold \(\tau_B > 0\) appears as an
active constraint at any sustainable equilibrium. The shadow price
structure would split into a scarcity component \(\lambda_j^*\) and a
flow-dependency component \(\pi_{ij}^*\), the latter measuring the
marginal downstream value of production. Such an extension would also
allow the model to engage directly with \citet{GeorgescuRoegen1971}'s
critique that neoclassical production functions ignore thermodynamic
constraints.

\textbf{Dynamic and stochastic environments.} The present results assume
static resource endowments. Extensions to time-varying resources
\(\mathbf{R}(t)\), stochastic scarcity, and Bayesian updating on
opponent types \citep{OsborneRubinstein1994} would test the robustness
of the variational equilibrium under distributional uncertainty. The
shadow price structure should extend to stochastic settings via robust
optimization or chance-constrained programming.

\textbf{Heterogeneous agent populations.} The two-agent worked examples
illustrate the general \(N\)-agent theorems but do not capture the full
complexity of large heterogeneous populations. Computational studies
with diverse agent types, varying metabolic efficiency, resource needs,
and constraint structures, would test the practical scalability of the
variational equilibrium.

\section{Conclusion}\label{conclusion}

We have shown that boundary constraints are not design choices but
necessary for multi-agent dissipative optimization. When multiple
self-maintaining open thermodynamic systems share a finite-resource
environment, no feasible joint allocation exists in which all agents
simultaneously achieve their unconstrained optima
(Theorem~\ref{thm-impossibility-infinite-freedom}). Each boundary
constraint carries a computable shadow price measuring its energetic
cost to the constrained agent (Theorem~\ref{thm-shadow-price}), and the
constrained system admits a variational equilibrium in which all agents
face the same shadow price for each shared resource
(Theorem~\ref{thm-variational-equilibrium}). This equilibrium is not
merely stable; it uniquely maximizes total social welfare over the
feasible set (Theorem~\ref{thm-welfare-optimality}). The welfare value
of the constraint system scales monotonically with scarcity
(Corollary~\ref{cor-constraint-scarcity-scaling}): under abundance,
constraints are costless; under scarcity, they become the dominant
determinant of system welfare. All results are denominated in energy
units, obtained entirely within the framework of constrained
optimization (KKT conditions, GNEP, and variational equilibrium theory),
and computationally verified.

These results establish the mathematical foundation for the subsequent
papers in this series. \emph{Thermodynamic Friction} \citep{TC-IV}
quantifies the energy cost of constraint violation. Strategic entropy
injection \citep{TC-II} formalizes the cost of information corruption.
Accumulated negentropy \citep{TC-III} establishes the irreplaceability
of complex dissipative structures. Cooperative equilibrium \citep{TC-V}
proves that cooperation is the unique efficient Nash Equilibrium under
energy-denominated payoffs. Value dynamics \citep{TC-VI} characterizes
stable coexistence as a dynamical attractor. The constraint necessity
and shadow price results derived here, the proof that boundary
constraints are unavoidable and that each carries a computable cost,
provide the structural foundation on which all subsequent results
depend.

\addcontentsline{toc}{section}{References}
\bibliography{../../shared/main}
\end{document}
